% ═══════════════════════════════════════════════════════════════
% LEHRERHINWEISE: Las preposiciones de lugar
% Spanisch Klasse 8 - Unidad 6, DS 2 (Einzelstunde, 45 Min)
% ═══════════════════════════════════════════════════════════════

\documentclass[11pt, a4paper]{article}

\usepackage[utf8]{inputenc}
\usepackage[T1]{fontenc}
\usepackage[ngerman]{babel}
\usepackage{helvet}
\renewcommand{\familydefault}{\sfdefault}

\usepackage[margin=2cm]{geometry}
\usepackage{enumitem}
\usepackage{tabularx}
\usepackage{booktabs}
\usepackage{longtable}

\usepackage{xcolor}
\usepackage{amssymb}
\usepackage{tcolorbox}
\tcbuselibrary{skins}

\usepackage{fancyhdr}
\usepackage{lastpage}
\usepackage{needspace}

% ═══════════════════════════════════════════════════════════════
% FARBEN
% ═══════════════════════════════════════════════════════════════

\definecolor{spanisch-color}{HTML}{c41e3a}
\definecolor{grau-color}{HTML}{888888}
\definecolor{hellgrau-color}{HTML}{f5f5f5}
\definecolor{basis-color}{HTML}{2a7a4b}
\definecolor{standard-color}{HTML}{2c5aa0}
\definecolor{erweiterung-color}{HTML}{7b1fa2}
\definecolor{loesung-color}{HTML}{c62828}

\colorlet{spanisch}{spanisch-color}
\colorlet{grau}{grau-color}
\colorlet{hellgrau}{hellgrau-color}
\colorlet{basis}{basis-color}
\colorlet{standard}{standard-color}
\colorlet{erweiterung}{erweiterung-color}
\colorlet{loesung}{loesung-color}

\newcommand{\fachfarbe}{spanisch}

% ═══════════════════════════════════════════════════════════════
% VARIABLEN
% ═══════════════════════════════════════════════════════════════

\newcommand{\thema}{Las preposiciones de lugar}
\newcommand{\sequenz}{06}
\newcommand{\block}{b}
\newcommand{\fach}{Spanisch}
\newcommand{\klasse}{8}
\newcommand{\dauer}{45 Min. (Einzelstunde)}
\newcommand{\lerngruppengroesse}{24}
\newcommand{\datum}{\today}

% ═══════════════════════════════════════════════════════════════
% KOPF- UND FUSSZEILE
% ═══════════════════════════════════════════════════════════════

\pagestyle{fancy}
\fancyhf{}
\renewcommand{\headrulewidth}{0.4pt}
\renewcommand{\footrulewidth}{0.4pt}

\lhead{\fach{} -- Klasse \klasse}
\chead{\textbf{LH-\sequenz\block{} (Lehrerhinweise)}}
\rhead{\datum}

\lfoot{}
\cfoot{}
\rfoot{Seite \thepage\ von \pageref{LastPage}}

% ═══════════════════════════════════════════════════════════════
% BEFEHLE
% ═══════════════════════════════════════════════════════════════

\newcommand{\B}{\textcolor{basis}{\textbf{[B]}}}
\newcommand{\Std}{\textcolor{standard}{\textbf{[S]}}}
\newcommand{\E}{\textcolor{erweiterung}{\textbf{[E]}}}
\newcommand{\lsg}[1]{\textcolor{loesung}{\textbf{#1}}}
\newcommand{\es}[1]{\textcolor{\fachfarbe}{\textbf{#1}}}

\newtcolorbox{kurzplan}{
    colback=hellgrau,
    colframe=\fachfarbe,
    colbacktitle=white,
    coltitle=\fachfarbe,
    boxrule=1pt,
    arc=0pt,
    fonttitle=\bfseries,
    attach boxed title to top left={yshift=-2mm, xshift=5mm},
    boxed title style={boxrule=0pt, colback=white},
    title=Kurzplan
}

\newtcolorbox{differenzierung}{
    colback=white,
    colframe=grau,
    colbacktitle=white,
    coltitle=grau,
    boxrule=0.5pt,
    arc=0pt,
    fonttitle=\bfseries,
    attach boxed title to top left={yshift=-2mm, xshift=5mm},
    boxed title style={boxrule=0pt, colback=white},
    title=Differenzierung
}

\newtcolorbox{fehlerbox}{
    colback=white,
    colframe=loesung,
    colbacktitle=white,
    coltitle=loesung,
    boxrule=0.5pt,
    arc=0pt,
    fonttitle=\bfseries,
    attach boxed title to top left={yshift=-2mm, xshift=5mm},
    boxed title style={boxrule=0pt, colback=white},
    title=Häufige Fehler
}

\newtcolorbox{materialbox}{
    colback=white,
    colframe=\fachfarbe,
    colbacktitle=white,
    coltitle=\fachfarbe,
    boxrule=0.5pt,
    arc=0pt,
    fonttitle=\bfseries,
    attach boxed title to top left={yshift=-2mm, xshift=5mm},
    boxed title style={boxrule=0pt, colback=white},
    title=Material-Checkliste
}

\newtcolorbox{spielebox}{
    colback=white,
    colframe=basis,
    colbacktitle=white,
    coltitle=basis,
    boxrule=0.5pt,
    arc=0pt,
    fonttitle=\bfseries,
    attach boxed title to top left={yshift=-2mm, xshift=5mm},
    boxed title style={boxrule=0pt, colback=white},
    title=Spielvorschläge (ad-hoc)
}

% ═══════════════════════════════════════════════════════════════
% DOKUMENT
% ═══════════════════════════════════════════════════════════════

\begin{document}

% ═══════════════════════════════════════════════════════════════
% KOPFBEREICH
% ═══════════════════════════════════════════════════════════════

\begin{center}
    {\Large\bfseries\textcolor{\fachfarbe}{Lehrerhinweise: \thema}}\\[0.3em]
    {\normalsize LH-\sequenz\block{} | \dauer{} | Qué pasa 1, S. 104--107}
\end{center}

\vspace{10pt}

% ═══════════════════════════════════════════════════════════════
% 1. KURZPLAN
% ═══════════════════════════════════════════════════════════════

\section*{1. Stundenverlauf (45 Min.)}

\begin{tabularx}{\textwidth}{|c|l|X|l|}
\hline
\textbf{Zeit} & \textbf{Phase} & \textbf{Inhalt} & \textbf{Material} \\
\hline
0--5 & Einstieg & Wiederholung Orte (Blitzrunde): L zeigt Bild $\rightarrow$ SuS nennen Ort & Bilder \\
\hline
5--15 & Input & L führt 10 Präpositionen mit Gegenstand ein (Stift + Box). SuS wiederholen. $\rightarrow$ \textbf{LB/AB: Tabelle} ausfüllen & LB, Stift, Box \\
\hline
15--22 & Übung 1 & \textbf{Buch S. 106, Ü2B:} 10 Bilder von Copito beschreiben. $\rightarrow$ \textbf{AB Aufg. 1} & Buch, AB \\
\hline
22--25 & Überleitung & L erklärt Aufgabe 4 (S. 107): Zwei Barrio-Bilder + Audio & Buch \\
\hline
25--32 & Hörverstehen & \textbf{Buch S. 107, Ü4:} Audio abspielen. SuS identifizieren Barrio A oder B & Audio \\
\hline
32--40 & Produktion & Anderes Barrio beschreiben (min. 5 Sätze). $\rightarrow$ \textbf{AB Aufg. 2} & AB \\
\hline
40--45 & Sicherung & 2-3 SuS lesen Beschreibung vor. Klasse prüft: Richtig? & -- \\
\hline
\end{tabularx}

% ═══════════════════════════════════════════════════════════════
% 2. DEUTSCHE ANWEISUNGEN FÜR BUCHÜBUNGEN
% ═══════════════════════════════════════════════════════════════

\section*{2. Deutsche Anweisungen für Buchübungen}

\textbf{Übung 2B (S. 106) -- Copito:}
\begin{quote}
\textit{Schau dir die 10 Bilder von Copito an. Beschreibe für jedes Bild, wo der Hund ist.}\\
\textbf{Beispiel:} Bild 1 $\rightarrow$ "Copito está debajo de la mesa."\\
Arbeitet zu zweit. Partner A beschreibt Bild 1-5, Partner B beschreibt Bild 6-10. Kontrolliert euch gegenseitig.
\end{quote}

\textbf{Übung 4 (S. 107) -- Dos barrios:}
\begin{quote}
\textbf{Schritt 1:} Hört das Audio. Welches Viertel wird beschrieben -- A oder B?\\
\textbf{Schritt 2:} Beschreibt danach das ANDERE Viertel (das nicht im Audio war) mit mindestens 5 Sätzen.\\
\textit{Redemittel:} El/La \_\_\_ está enfrente de... / al lado de... / entre... y...
\end{quote}

% ═══════════════════════════════════════════════════════════════
% 3. DIFFERENZIERUNG
% ═══════════════════════════════════════════════════════════════

\section*{3. Differenzierung}

\begin{tabularx}{\textwidth}{|l|X|X|X|}
\hline
& \B{} \textbf{Basis} & \Std{} \textbf{Standard} & \E{} \textbf{Erweiterung} \\
\hline
\textbf{AB Aufg. 1} & Nur 5 Bilder (1-5), Präpositionen vorgegeben & Alle 10 Bilder & + Wo ist Copito NICHT? (Verneinung) \\
\hline
\textbf{AB Aufg. 2} & 3 Sätze, Satzanfänge vorgegeben & 5 Sätze & 7+ Sätze, komplexere Strukturen \\
\hline
\end{tabularx}

% ═══════════════════════════════════════════════════════════════
% 4. HÄUFIGE FEHLER
% ═══════════════════════════════════════════════════════════════

\section*{4. Häufige Fehler}

\begin{fehlerbox}
\begin{tabularx}{\textwidth}{|l|X|X|}
\hline
\textbf{Fehler} & \textbf{Beispiel} & \textbf{Reaktion} \\
\hline
*de el & *"al lado de el parque" statt \es{al lado del parque} & Regel: de + el = del \\
\hline
*a el & *"enfrente a el cine" statt \es{enfrente del cine} & Regel: a + el = al, de + el = del \\
\hline
Falsche Präp. & *"encima la mesa" statt \es{encima de la mesa} & Fast alle brauchen "de"! \\
\hline
Links/Rechts & Verwechslung izquierda/derecha & Mit Händen zeigen \\
\hline
\end{tabularx}
\end{fehlerbox}

% ═══════════════════════════════════════════════════════════════
% 5. SPIELVORSCHLÄGE
% ═══════════════════════════════════════════════════════════════

\section*{5. Spielvorschläge (ad-hoc, Partnerarbeit)}

\begin{spielebox}
\textbf{1. "¿Dónde está?" Rätsel (5-8 Min)}
\begin{itemize}
    \item S1 versteckt gedanklich einen Gegenstand im Klassenzimmer
    \item S2 fragt: "¿Está debajo de la mesa?" -- "No."
    \item S2: "¿Está encima del armario?" -- "¡Sí!"
    \item Max. 5 Fragen, dann Rollentausch
\end{itemize}

\vspace{0.5em}

\textbf{2. Schnelles Diktat mit Bewegung (5 Min)}
\begin{itemize}
    \item L ruft Präposition $\rightarrow$ SuS zeigen mit Händen die Position
    \item "encima de" $\rightarrow$ Hände nach oben
    \item "debajo de" $\rightarrow$ Hände nach unten
    \item Wer falsch zeigt, setzt sich. Letzte 5 Stehende gewinnen.
\end{itemize}

\vspace{0.5em}

\textbf{3. Partner-Zeichnen (8-10 Min)}
\begin{itemize}
    \item S1 hat ein einfaches Barrio-Bild (5-6 Orte)
    \item S1 beschreibt: "La biblioteca está entre el parque y el cine."
    \item S2 zeichnet nach Beschreibung (ohne zu sehen)
    \item Vergleich am Ende: Wie ähnlich sind die Bilder?
\end{itemize}
\end{spielebox}

% ═══════════════════════════════════════════════════════════════
% 6. MATERIAL-CHECKLISTE
% ═══════════════════════════════════════════════════════════════

\section*{6. Material-Checkliste}

\begin{materialbox}
\begin{itemize}[label=$\square$]
    \item LB-06b-preposiciones.pdf (\lerngruppengroesse{} Kopien)
    \item AB-06b-preposiciones.pdf (\lerngruppengroesse{} Kopien)
    \item ML-06b-preposiciones.pdf (1× für Lehrkraft)
    \item Lehrwerk: Qué pasa 1, S. 104--107
    \item Audio-Track zu Übung 4 (S. 107)
    \item Gegenstand zur Demonstration (Stift, Box)
    \item Optional: PPT-06b-preposiciones (Beamer)
\end{itemize}
\end{materialbox}

% ═══════════════════════════════════════════════════════════════
% 7. LÖSUNGEN
% ═══════════════════════════════════════════════════════════════

\section*{7. Lösungen AB-\sequenz\block}

\textbf{Präpositionentabelle:}

\begin{tabular}{ll|ll}
\es{enfrente de} & \lsg{gegenüber von} & \es{a la derecha de} & \lsg{rechts von} \\
\es{encima de} & \lsg{über, auf} & \es{entre} & \lsg{zwischen} \\
\es{al lado de} & \lsg{neben} & \es{detrás de} & \lsg{hinter} \\
\es{a la izquierda de} & \lsg{links von} & \es{debajo de} & \lsg{unter} \\
\es{en} & \lsg{in, an, auf} & \es{delante de} & \lsg{vor} \\
\end{tabular}

\vspace{1em}

\textbf{Aufgabe 1: Copito} (10P) -- Beispiele (abhängig von Buchbildern):
\begin{itemize}
    \item Bild 1: \lsg{Copito está debajo de la mesa.}
    \item Bild 2: \lsg{Copito está encima de la silla.}
    \item Bild 3: \lsg{Copito está delante de la puerta.}
    \item usw. (1P pro korrektem Satz)
\end{itemize}

\vspace{0.5em}

\textbf{Aufgabe 2: Barrio} (10P) -- Mögliche Sätze:
\begin{itemize}
    \item \lsg{La iglesia está enfrente de la plaza.}
    \item \lsg{El supermercado está al lado de la farmacia.}
    \item \lsg{La biblioteca está a la derecha del parque.}
    \item \lsg{El cine está entre el restaurante y la cafetería.}
    \item \lsg{La panadería está a la izquierda del banco.}
\end{itemize}

(2P pro Satz: 1P Struktur, 1P passende Präposition)

\end{document}
